% Generated draft: requires mathematical and source review.
\documentclass[11pt]{article}
\usepackage[margin=25mm]{geometry}
\usepackage{fontspec}
\setmainfont{Latin Modern Roman}
\usepackage{amsmath,amssymb,mathtools,unicode-math}
\setmathfont{Latin Modern Math}
\usepackage{enumitem,needspace,xurl,hyperref}
\setlength{\parindent}{0pt}
\setlength{\parskip}{5pt}
\setlength{\emergencystretch}{3em}
\newcommand{\N}{\mathbb{N}}
\newcommand{\Z}{\mathbb{Z}}
\newcommand{\Q}{\mathbb{Q}}
\newcommand{\R}{\mathbb{R}}
\newcommand{\C}{\mathbb{C}}
\newcommand{\F}{\mathbb{F}}
\newcommand{\A}{\mathbb{A}}
\newcommand{\PP}{\mathbb{P}}
\newcommand{\E}{\mathbb{E}}
\newcommand{\eps}{\varepsilon}
\newcommand{\dd}{\,\mathrm{d}}
\begin{document}
\section*{UnsolvedMath ID: 30003727}
\textbf{Source problem code:} catalogue entry
\subsection{Definitions and mathematical statement}
Let \(G_N\in\C^{N\times N}\) be a complex Ginibre matrix, meaning a matrix with independent, identically distributed centered circular complex Gaussian entries. The entry variance, including its dependence on \(N\), is not specified in the supplied record. Write its eigenvalues as \(\lambda_1,\ldots,\lambda_N\). For each eigenvalue under consideration, let \(r_i,l_i\in\C^N\setminus\{0\}\) be right and left eigenvectors, respectively, with
\[
G_Nr_i=\lambda_i r_i,
\qquad l_i^*G_N=\lambda_i l_i^*,
\]
where \({}^*\) denotes conjugate transpose. These equations do not fix the scaling of either eigenvector.

The supplied problem asks for asymptotics, as \(N\to\infty\), of the volume of a parallelepiped formed by eigenvectors associated with eigenvalues near fixed deterministic parameters \(\nu_1,\ldots,\nu_\ell\in\C\), and for the dependence on \(N\) of the minimum angle between eigenvectors. The record does not impose distinctness or a spectral location condition on the parameters.

To describe possible observables precisely without selecting an unsupported interpretation, let \(J\subseteq\{1,\ldots,N\}\) be a specified set of indices, and choose either \(v_i=r_i\) for every \(i\in J\), or \(v_i=l_i\) for every \(i\in J\). If unit normalization and the Hermitian Gram-volume convention are intended, set
\[
 u_i=\frac{v_i}{\|v_i\|_2},\qquad
 V_N(J)=\left(\det\bigl[u_i^*u_j\bigr]_{i,j\in J}\right)^{1/2}.
\]
If the angle is intended to be the phase-invariant angle between complex lines, then, for \(|J|\geq2\), its minimum is
\[
 \theta_N(J)=\min_{\substack{i,j\in J\\i\ne j}}
 \arccos\bigl(|u_i^*u_j|\bigr)\in[0,\pi/2].
\]
These are conditional definitions of natural candidate observables; the supplied record does not establish that they are the intended conventions. In particular, an unnormalized volume depends on the arbitrary eigenvector scalings, and a Euclidean real-volume interpretation can give a different observable.

\textbf{Statement requiring specification.} Determine the large-\(N\) behavior of the intended eigenvector volume for eigenvalues near \(\nu_1,\ldots,\nu_\ell\), and determine the large-\(N\) behavior of the intended minimum eigenvector angle. A unique precise asymptotic assertion cannot be extracted until the matrix scaling, neighborhood sizes, eigenvalue selection or conditioning rule, eigenvector normalization, volume and angle conventions, participating left/right families, and mode of asymptotic comparison are specified. The minimum may concern the selected eigenvectors or the entire spectrum; the supplied text does not decide between these possibilities.
\subsection{Short English statement}
The problem asks how eigenvectors of a large complex Ginibre matrix span volume near prescribed spectral locations, and how small their mutual angles become. Its formulation leaves essential normalization, selection, geometric, and asymptotic conventions unresolved. The displayed observables describe possible conventions only. This is a statement-only draft requiring human mathematical review.
\subsection{Sources}
\begin{itemize}
\item Supplied local UnsolvedMath registry record, permanent ID \(30003727\), title ``Eigenvector Volumes and Angles in Ginibre Matrices,'' exact source code \texttt{catalogue entry}, in the Oberwolfach Reports Open Problems collection. No individual proposer is identified; the supplied proposed year is 2018.
\item Supplied source citation: \emph{Strongly Correlated Random Interacting Processes} (2019), Oberwolfach Reports, pp.\ 187--253, DOI \texttt{10.4171/owr/2018/4}. Supplied source and background link: \url{https://doi.org/10.4171/owr/2018/4}. The background describes Oberwolfach Reports as documenting mathematical workshops and their open-problem discussions.
\item Supplied literature source labeled ``original OWR report'': \url{https://doi.org/10.4171/OWR/2018/4}.
\item Supplied literature source labeled ``Crawford--Rosenthal, Ginibre eigenvector correlations'': \url{https://arxiv.org/abs/1805.08993}.
\item Literature assessment dated exactly \texttt{2026-08-21}: the supplied curation labels the problem ``partially solved'' and states that Crawford--Rosenthal derive broad asymptotics for Ginibre eigenvector correlations, while the specific selected-eigenvector volume and minimal-angle asymptotics asked here are not completely determined. This is supplied dated curation, not an independently verified literature assessment; the source explicitly requests independent verification of time-sensitive status and cited claims.
\item The supplied statement status is \texttt{corrected\_verified}, accompanied by a claimed manual semantic, typographic, and LaTeX review asserting self-containment and preservation of the mathematical target. That supplied review label does not resolve the formulation issues recorded here or constitute mathematical verification of this draft.
\end{itemize}
\end{document}
